\documentclass[11pt]{tex/lawoutline}
\usepackage{amssymb}
\usepackage{enumitem}

\course{Civil Procedure}
\student{Prof: J. Greiner}
\semester{Fall 2026}

\begin{document}

\maketitle
\tableofcontents


\question
{\textit{J. McIntyre}, Foreign Companies}
{
    Imagine that you are general counsel for a foreign company that manufactures dangerous products. The company's CEO wants to sell as many products as possible to customers throughout the United States but does not want the company to be sued in any American state court.

    \begin{itemize}
    \item \textbf{Does the plurality provide a simple way to avoid jurisdiction?}
    \item Not completely, but the company could reduce its risk by avoiding direct sales, advertising, or other activities specifically aimed at individual states.

    \item \textbf{Could the company target the United States generally while avoiding specific states?}
    \item Under the plurality's reasoning, probably. Targeting the United States as a whole does not necessarily mean purposefully targeting every individual state.

    \item \textbf{Could the company still be sued in some states?}
    \item Yes. A state might exercise jurisdiction using a \textbf{long-arm statute}, or if the company directly targeted that state, regularly sold products there, maintained contacts there, consented to jurisdiction, etc.

    \item \textbf{Which branch of government could change this result?}
    \item Congress or state legislators could enact legislation or long-arm statutes that authorize jurisdiction over defendants.

    \item \textbf{What would that mean?}
    \item It would be easier for a state court to hear the case even when the defendant lacked sufficient contacts with the particular state where the court was located.
\end{itemize}
}

\end{document}